% Generated draft: requires mathematical and source review.
\documentclass[11pt]{article}
\usepackage[margin=25mm]{geometry}
\usepackage{fontspec}
\setmainfont{Latin Modern Roman}
\usepackage{amsmath,amssymb,mathtools,unicode-math}
\setmathfont{Latin Modern Math}
\usepackage{enumitem,needspace,xurl,hyperref}
\setlength{\parindent}{0pt}
\setlength{\parskip}{5pt}
\setlength{\emergencystretch}{3em}
\newcommand{\N}{\mathbb{N}}
\newcommand{\Z}{\mathbb{Z}}
\newcommand{\Q}{\mathbb{Q}}
\newcommand{\R}{\mathbb{R}}
\newcommand{\C}{\mathbb{C}}
\newcommand{\F}{\mathbb{F}}
\newcommand{\A}{\mathbb{A}}
\newcommand{\PP}{\mathbb{P}}
\newcommand{\E}{\mathbb{E}}
\newcommand{\eps}{\varepsilon}
\newcommand{\dd}{\,\mathrm{d}}
\begin{document}
\section*{UnsolvedMath ID: 2521}
\textbf{Source problem code:} catalogue entry
\subsection{Definitions and mathematical statement}
Let \(\omega=\{0,1,2,\ldots\}\), and let \(U\) be a nonprincipal ultrafilter on \(\omega\): it is closed under finite intersections and supersets, contains no finite set, and for every \(X\subseteq\omega\) contains exactly one of \(X\) and its complement. For groups \((G_i)_{i\in\omega}\), define
\[
N_U=\{(g_i)\in\prod_{i\in\omega}G_i:\{i:g_i=1_{G_i}\}\in U\},
\qquad
\prod_U G_i=\left(\prod_{i\in\omega}G_i\right)/N_U.
\]
The notation \(\Z^\omega/U\) denotes the corresponding ultrapower of the additive group \(\Z\). A homomorphic image means the target of a surjective group homomorphism, with no topology imposed.

The question, attributed to G. M. Bergman, is whether for every such \(U\) and every group \(B\),
\[
\left[\forall b\in B\ \exists H_b\leq B\ \exists\pi_b:\Z^\omega/U\twoheadrightarrow H_b,\quad b\in H_b\right]
\]
implies
\[
\exists(G_i)_{i\in\omega}\ \exists\Pi:\prod_U G_i\twoheadrightarrow B.
\]
The subgroups and maps may depend on \(b\); the ultrafilter in the desired global representation is the same fixed \(U\) used in the premise.

The source records a negative answer by Corson (2025). Its supplied literature triage dated 2026-08-17 reports the torsion counterexample
\[
B=\bigoplus_{n\geq1}\Z/n\Z,
\]
which satisfies the local condition but is not a homomorphic image of a nonprincipal ultraproduct over \(\omega\), uniformly over the choice of nonprincipal ultrafilter. It notes that Corson's restatement permits a different ultrafilter for a putative global representation, but the uniform non-representability conclusion is stronger than needed for this same-\(U\) question. The scope here remains the countable index set \(\omega\), not arbitrary ultrafilter index sets.

The background identifies the earlier question as Bergman, Question 19 (2015), and says an affirmative answer would have represented every torsion group for every \(U\), thereby answering both parts of Kourovka 21.11 affirmatively. That is preserved only as a historical conditional implication, not an affirmative outcome. Both the registry and the dated triage classify the stated question solved negatively. The draft reports the supplied published counterexample without independently checking its proof.
\subsection{Short English statement}
If every element of a group lies in a subgroup which is an image of one fixed integer ultrapower, must the entire group be an image of a group ultraproduct using the same ultrafilter? The supplied Corson result answers no, with a torsion counterexample excluded uniformly for nonprincipal ultrafilters on \(\omega\).
\subsection{Sources}
\begin{itemize}
\item Local UnsolvedMath registry, permanent ID 2521, source code catalogue entry, ``Kourovka Notebook Problem 21.12,'' drawn from \emph{Kourovka Notebook}, New Problems, 21st issue (2026), original problem 21.12, attributed in the background to G. M. Bergman. The related reference to both parts of 21.11 has no separate supplied statement.
\item G. M. Bergman, \emph{Pacific Journal of Mathematics} 274 (2015), 451--495, Question 19; the source supplies no paper title or URL.
\item Samuel M. Corson, ``On homomorphic images of ultraproducts,'' \emph{Journal of Algebra} 681 (2025), 306--317: \url{https://doi.org/10.1016/j.jalgebra.2025.06.001}.
\item Samuel M. Corson, ``On homomorphic images of ultraproducts,'' author manuscript, arXiv:2503.09228 (2025): \url{https://arxiv.org/abs/2503.09228}.
\end{itemize}
\end{document}
